% !TeX root = ../BTechProjectTemplate.tex
\chapter{Implementation}

\section{Introduction}
This chapter describes the experimental setup, including dataset preparation, the preprocessing pipeline, and the multi-objective optimization strategy used to train the Swin-MMHCA model.

\section{Dataset and Subject-Level Partitioning}
The foundation of our research is the \textbf{IXI Dataset} (Information eXtraction from Images), which consists of nearly 600 MR images from healthy subjects collected at three prominent London hospitals: Hammersmith Hospital, Guy's Hospital, and the Institute of Psychiatry. Each subject underwent T1, T2, and PD-weighted scans using 1.5T and 3T Philips Medical Systems scanners.

The inclusion of multiple modalities is critical for super-resolution. T1-weighted scans, acquired using a modified driven equilibrium Fourier transform (MDEFT) sequence, provide high anatomical contrast between gray and white matter. T2-weighted scans, acquired using a dual-echo turbo spin-echo (TSE) sequence, are sensitive to fluid and pathological changes. By utilizing the IXI dataset, we ensure that our model learns from high-quality, research-grade clinical data that captures the natural variability of the human brain's anatomy across a diverse population.

To prevent "data leakage" and ensure the generalizability of our Swin-MMHCA model, we implemented a rigorous \textbf{Subject-Level Split}. In many deep learning studies, slices from the same subject are mixed between training and testing sets, leading to biased results because the model "memorizes" the specific anatomy of those subjects. Our partition ensures that a subject's unique neuro-anatomical features are seen in only one phase:
\begin{itemize}
    \item \textbf{Training Set (80\%):} Used to optimize the weights of the SwinV2 backbone and the MMHCA module, exposing the model to a wide variety of cortical patterns.
    \item \textbf{Validation Set (10\%):} Used during training for early stopping and hyperparameter tuning to ensure the model does not overfit to the training subjects.
    \item \textbf{Test Set (10\%):} Reserved for the final, unbiased evaluation of the model's performance on completely unseen anatomies, simulating a real-world clinical deployment.
\end{itemize}

\section{SimpleITK Preprocessing and Registration Details}
The SimpleITK library provided the core functionality for our research-grade preprocessing pipeline. This stage is critical because any misalignment between modalities would lead to the propagation of artifacts through the MMHCA module.

\subsection{Rigid Registration and Mutual Information}
We utilized the \textit{Mattes Mutual Information} metric to drive the registration process. Mutual Information is an entropy-based metric that measures the statistical dependence between two images. Unlike Mean Squared Error, which requires similar intensity values, Mutual Information can align images with completely different contrasts (e.g., T1 and T2).
The transformation model was set to \textit{AffineTransform(3)} to perform robust 3D registration. This ensures that the T1 and PD modalities are perfectly aligned with the T2 volume in three-dimensional space before any slicing occurs, maintaining anatomical consistency across the entire brain. We used the \textit{RegularStepGradientDescentOptimizerv4} to iteratively refine the transformation parameters.

\subsection{2D Slice Extraction and Preprocessing Filters}
Following the 3D registration, we converted and extracted 2D slices from the aligned volumes for training and evaluation. Each slice underwent a series of filtering operations:
\begin{enumerate}
    \item \textbf{Intensity Rescaling:} We used the \textit{RescaleIntensityImageFilter} to map the raw 12-bit or 16-bit MRI values into a standard float $[0, 1]$ range.
    \item \textbf{Gaussian Smoothing:} A slight Gaussian kernel ($\sigma = 0.5$) was optionally applied during registration to smooth the cost function and avoid local minima.
    \item \textbf{Brain Masking:} We utilized the \textit{BinaryThresholdImageFilter} to generate a rough brain mask, which was then used to reject slices containing less than 28\% brain tissue.
\end{enumerate}
This rigorous preprocessing ensures that the Swin-MMHCA model is trained on high-quality, perfectly aligned anatomical features.

\section{Training Protocol and Objective Functions}
The model is trained using a multi-objective optimization strategy. The total loss function $\mathcal{L}_{total}$ is:
\begin{equation}
    \mathcal{L}_{total} = \lambda_1 \mathcal{L}_{L1} + \lambda_2 \mathcal{L}_{edge} + \lambda_3 \mathcal{L}_{perc} + \lambda_4 \mathcal{L}_{seg} + \lambda_5 \mathcal{L}_{det}
\end{equation}

\subsection{Structural Fidelity: Charbonnier L1 Loss}
We utilize the Charbonnier loss, a differentiable approximation of $L_1$:
\begin{equation}
    \mathcal{L}_{L1}(\hat{X}, X) = \frac{1}{CHW} \sum_{c,h,w} \sqrt{(\hat{X}_{c,h,w} - X_{c,h,w})^2 + \epsilon^2}
\end{equation}
It is more robust to sharp intensity transitions than $L_2$ loss.

\subsection{High-Frequency Edge Preservation Loss}
To preserve anatomical boundaries, we compute the gradient maps using Sobel operators and minimize the $L_1$ distance in the gradient domain:
\begin{equation}
    \mathcal{L}_{edge} = \mathbb{E}[\|\nabla \hat{X} - \nabla X\|_1]
\end{equation}

\subsection{Perceptual Fidelity: LPIPS Loss}
We incorporate the Learned Perceptual Image Patch Similarity (LPIPS) loss, which computes the distance between images in a deep feature space, matching human perceptual judgments.

\subsection{Anatomical Semantic Loss: Dice-BCE}
For the segmentation auxiliary task, we utilize a combination of Binary Cross Entropy (BCE) and the Dice coefficient:
\begin{equation}
    \mathcal{L}_{seg} = \mathcal{L}_{BCE} + \mathcal{L}_{Dice}
\end{equation}

\subsection{Pathological Awareness: Focal Loss}
Detecting high-variance regions (lesions) is an imbalanced task. We utilize Focal Loss to prevent the model from being dominated by "easy" background pixels:
\begin{equation}
    \mathcal{L}_{focal} = -\alpha (1 - p_t)^\gamma \log(p_t)
\end{equation}

\section{Evaluation Metrics and Mathematical Definitions}
To rigorously assess the performance of the super-resolution model, we utilize three complementary metrics that evaluate different aspects of image fidelity: pixel-wise accuracy, structural similarity, and perceptual quality.

\subsection{Peak Signal-to-Noise Ratio (PSNR)}
PSNR is the most common metric used to measure the quality of reconstruction in lossy compression and image restoration. It is defined based on the Mean Squared Error (MSE) between the ground truth image ($X$) and the super-resolved image ($\hat{X}$):
\begin{equation}
    MSE = \frac{1}{mn} \sum_{i=0}^{m-1} \sum_{j=0}^{n-1} [X(i,j) - \hat{X}(i,j)]^2
\end{equation}
\begin{equation}
    PSNR = 10 \cdot \log_{10} \left( \frac{MAX_I^2}{MSE} \right)
\end{equation}
where $MAX_I$ is the maximum possible pixel value (e.g., 1.0 for normalized images). While PSNR is effective for measuring signal-to-noise levels, it does not always align with human visual perception.

\subsection{Structural Similarity Index (SSIM)}
SSIM is a perception-based model that considers image degradation as perceived change in structural information. Unlike MSE, which measures absolute error, SSIM compares three components: luminance ($l$), contrast ($c$), and structure ($s$):
\begin{equation}
    SSIM(x, y) = \frac{(2\mu_x\mu_y + C_1)(2\sigma_{xy} + C_2)}{(\mu_x^2 + \mu_y^2 + C_1)(\sigma_x^2 + \sigma_y^2 + C_2)}
\end{equation}
where $\mu$ is the mean, $\sigma^2$ is the variance, and $C_1, C_2$ are constants to stabilize the division. SSIM ranges from -1 to 1, with 1 indicating perfect identity.

\subsection{Learned Perceptual Image Patch Similarity (LPIPS)}
LPIPS measures the distance between two images in a deep feature space. It utilizes a pre-trained network (e.g., AlexNet or VGG) to extract activations from multiple layers. The distance is calculated as a weighted average of the $L_2$ differences between normalized features:
\begin{equation}
    d(x, x_0) = \sum_l \frac{1}{H_l W_l} \sum_{h,w} \| w_l \odot (\hat{y}_{hw}^l - \hat{y}_{0hw}^l) \|_2^2
\end{equation}
A lower LPIPS value indicates that the images are perceptually similar to a human observer, making it ideal for evaluating the "look and feel" of medical textures.
The training configurations and architectural hyperparameters are critical for the stable convergence of the SwinV2 backbone. We utilized an initial warm-up period of 5 epochs followed by a cosine annealing learning rate schedule to fine-tune the weights. The detailed parameters are summarized in Table \ref{tab:hyperparameters}.

\begin{table}[H]
    \centering
    \caption{Training and Model Hyperparameters.}
    \label{tab:hyperparameters}
    \begin{tabular}{ll}
        \toprule
        \textbf{Hyperparameter} & \textbf{Value} \\
        \midrule
        Optimizer & AdamW \\
        Base Learning Rate & $1 \times 10^{-4}$ \\
        Weight Decay & $1 \times 10^{-2}$ \\
        Batch Size & 16 \\
        Total Epochs & 250 \\
        Patch Size & $64 \times 64$ \\
        Embedding Channels & 192 \\
        Swin Blocks per Stage & [2, 2, 6, 2] \\
        Attention Heads & [3, 6, 12, 24] \\
        Window Size & 8 \\
        Loss Weight ($\lambda_{L1}$) & 1.0 \\
        Loss Weight ($\lambda_{edge}$) & 0.1 \\
        Loss Weight ($\lambda_{seg}$) & 0.5 \\
        \bottomrule
    \end{tabular}
\end{table}
The use of the AdamW optimizer with weight decay helped regularize the large parameter space of the hierarchical transformer, preventing overfitting to the specific textures of the training subjects.
